
In recent years, sentiment analysis in social media has attracted a lot of research interest and has been used for a number of applications. Unfortunately, research has been hindered by the lack of suitable datasets, complicating the comparison between approaches. To address this issue, we have proposed \emph{SemEval-2013 Task 2: Sentiment Analysis in Twitter},
%, which introduces a new dataset and thus allows to compare different methods on the same data.
which included two subtasks: A, an expression-level subtask, and B, a message-level subtask. We used crowdsourcing on Amazon Mechanical Turk to label a large Twitter training dataset along with additional test sets of Twitter and SMS messages for both subtasks. All datasets used in the evaluation are released to the research community.
The task attracted significant interest and a total of 149 submissions from 44 teams. The best-performing team achieved an F1 of 88.9\% and 69\% for subtasks A and B, respectively.

%In recent years it has become evident that sentiment analysis in social media is very useful. Although this research problem has been explored, the authors are unaware of a formal evaluation to compare how well different sentiment system achieve this task. In this paper, we describe the definition, data, evaluation, and results for
% the inaugural \emph{SemEval-2013 Task 2: Sentiment Analysis in Twitter} Task. The task included two subtasks; an expression level and message level task.
%% The task included two sub-tasks; an expression-level task, targeting contextual polarity disambiguation where
%% participants were given a message and a marked instance of a word or phrase, and were asked to determine
%% whether it is positive, negative or neutral in the context of the message. On the other hand, Task B, was a message polarity classification task where participants were given a message and were asked to say whether it contained positive, negative, or neutral sentiment.
%We used crowdsourcing on Amazon Mechanical Turk to label a large Twitter training dataset along with an additional test set of Twitter and SMS data for both tasks. The task was very successful, bringing in a total of 149 submissions were from 44 teams with the highest team results at 88.9\% and 69\% F-measure for Task A and B respectively.

%We used crowdsourcing on Amazon Mechanical Turk to develop two Twitter datasets and two additional SMS datasets (for task A and B). Participants were allowed to submit constrained (using only the provided training data) or unconstrained systems (using additional annotated training data). We scored the systems using an averaged F1-measure over the positive and the negative examples. The Task was very successful, bringing in A total of 149 submissions were from 44 teams with the highest team results at 88.9\% and 69\% F-measure for Task A and B respectively.
